\paragraph{Update design and attribution.}
The experimental comparison evaluates Block3 as an integrated block-level update: shared advantages, joint importance ratios and valid-length-weighted block reduction. Section~\ref{sec:method} characterizes the resulting cross-token credit, gradient scale and clipping geometry. Accordingly, the measured performance differences concern the complete update rather than advantage averaging alone. Component-specific attribution can be examined separately by holding ratios and normalization fixed when varying feedback sharing. The signal-estimation analysis describes averaging under explicit assumptions; task-level accuracy is measured by the benchmark evaluations.

\paragraph{Statistical unit and checkpoint reporting.}
Each reported arm uses one training seed. Paired-question intervals characterize evaluation uncertainty conditional on the trained checkpoints, rather than variability across training seeds. They are pointwise and use no multiple-comparison correction. Each AIME set contains 30 questions, so one additional solved question changes Pass@8 by 3.33 points. We retain the checkpoint trajectories alongside Step200, a common scheduled endpoint adopted retrospectively for reporting rather than a preregistered selection rule. Window and weighting variants arose through exploratory experiments; independent runs would evaluate the reproducibility of their observed differences.

\paragraph{Protocols and task coverage.}
Historical results are organized by their original prompt, seed-handling, teacher and implementation settings. Each stratum supplies a method comparison under its own protocol. The two fully retained comparisons share a DAPO-derived training pool; the overlap audit checks the retained data rather than semantic equivalence or upstream model exposure. The 4B completion study compares the two trained methods without a matched original-student evaluation, while the official Instruct pair additionally includes that control. Independent datasets, training seeds and teacher--student pairs provide natural axes for extending this evaluation.

\paragraph{Complementary diagnostics.}
Prompt probes examine behavior before optimization, while retained training trajectories characterize subsequent evolution. Entropy maps describe the prefixes visited by the current policy; missing boxed markers measure one specific formatting property. Gradient norms and length-stop rates track numerical and generation behavior, and sign disagreement measures how feedback is redistributed. Benchmark accuracy provides a separate task-level assessment. Interpreting these measurements jointly distinguishes input compatibility, optimization dynamics and mathematical correctness without treating one as a proxy for all three.

\paragraph{Baseline and computational scope.}
Sampled-token OPD is the matched experimental baseline; other OPD recipes are discussed as related work. Block3 changes how the available teacher feedback drives the update and still scores every valid response position. The present evaluation therefore concerns the behavior and accuracy of this update, not a reduction in teacher-scoring work. Per-task, per-checkpoint reporting captures the observed variation. Further comparisons can examine component choices and accuracy--cost trade-offs under a common budget.
